%% ============================================================
%% Rozdział 4 — Podsumowanie i wnioski
%% ============================================================
\chapter{Podsumowanie}\label{chapter:4_Summary}
W niniejszym rozdziale należy zamieścić syntetyczne podsumowanie całej pracy, obejmujące najważniejsze osiągnięcia, rezultaty i wnioski płynące z przeprowadzonych badań lub realizacji projektu. Celem tej części jest przedstawienie, w jakim stopniu zrealizowano założone cele oraz jakie znaczenie praktyczne lub naukowe mają uzyskane wyniki.

Podsumowanie powinno również obejmować refleksję nad napotkanymi trudnościami i ograniczeniami, które mogły mieć wpływ na przebieg prac lub uzyskane rezultaty. Wskazane jest również sformułowanie sugestii dotyczących możliwych kierunków dalszych badań, rozwoju projektu lub jego udoskonalenia.

\section{Wnioski}
Wnioski stanowią syntetyczne zestawienie kluczowych obserwacji i rezultatów uzyskanych w pracy. Powinny być przedstawione w formie uporządkowanej listy, z naciskiem na praktyczne znaczenie wyników oraz ich zgodność z postawionymi wcześniej celami i hipotezami.

\section{Kierunek dalszego rozwoju}
\noindent W tym podrozdziale należy wskazać możliwe kierunki rozbudowy, optymalizacji lub zastosowania opracowanego rozwiązania. Można tu także opisać:
\begin{itemize}
    \item potencjalne usprawnienia funkcjonalne lub techniczne
    \item możliwość integracji z innymi systemami
    \item zastosowanie wyników w praktyce inżynierskiej, naukowej lub przemysłowej
    \item dalsze badania lub testy umożliwiające weryfikację skuteczności opracowanego podejścia
\end{itemize}

\noindent Podsumowując, rozdział ten powinien stanowić końcowy, logiczny etap pracy, prezentujący jej dorobek oraz wskazujący, w jaki sposób może on zostać rozwinięty lub wykorzystany w przyszłości.
